% Mathématiques 4e — cours V3
% Dépendance entre grandeurs, tableaux et graphiques

\section{Objectifs}

\begin{itemize}
\item Reconnaître qu'une grandeur peut dépendre d'une autre.
\item Utiliser l'expression « en fonction de ».
\item Traduire une dépendance par une formule simple.
\item Produire et lire un tableau de valeurs.
\item Passer d'une formule à un tableau.
\item Placer des points issus d'un tableau dans un repère.
\item Lire une valeur sur une courbe ou un nuage de points.
\item Retrouver une valeur d'entrée correspondant à une valeur de sortie.
\item Comparer différentes représentations d'une même dépendance.
\item Distinguer une dépendance proportionnelle d'une dépendance non proportionnelle.
\end{itemize}

\section{1. Une grandeur dépend d'une autre}

Dans de nombreuses situations, la valeur d'une grandeur est déterminée par une autre.

Exemples :
\begin{itemize}
\item l'aire d'un carré dépend de son côté ;
\item le prix dépend de la quantité achetée ;
\item une température dépend de l'heure du relevé ;
\item la hauteur d'eau dépend du temps de remplissage.
\end{itemize}

On dit que la seconde grandeur est donnée \textbf{en fonction de} la première.

\section{2. Une règle de calcul}

Un service facture :
\[
3\ \text{€}
\]
par heure plus :
\[
5\ \text{€}
\]
de frais fixes.

Pour une durée $t$ :
\[
P=3t+5.
\]

La formule donne le prix en fonction de la durée.

\section{3. Plusieurs représentations}

Une même dépendance peut être décrite par :
\begin{itemize}
\item une phrase ;
\item une formule ;
\item un tableau de valeurs ;
\item un graphique.
\end{itemize}

\coursefigure{/images/mathematiques/4eme/fonctions-representations.svg}{Formule, tableau et graphique reliés entre eux.}{Changer de représentation ne change pas la relation entre les grandeurs ; chaque forme met en évidence des informations différentes.}

\section{4. Construire un tableau}

Avec :
\[
y=2x+3,
\]
on choisit des valeurs de $x$.

\begin{tabular}{c|c|c|c|c}
$x$ & $0$ & $1$ & $2$ & $5$ \\
$y$ & $3$ & $5$ & $7$ & $13$ \\
\end{tabular}

Chaque colonne donne un couple :
\[
(x;y).
\]

\section{5. Calculer une valeur}

Pour :
\[
x=4,
\]
dans :
\[
y=2x+3,
\]
on obtient :
\[
y=2\times4+3=11.
\]

La valeur de sortie correspondant à :
\[
4
\]
est donc :
\[
11.
\]

\section{6. Retrouver une entrée}

On cherche la valeur de $x$ qui produit :
\[
y=17
\]
dans :
\[
y=2x+3.
\]

On résout :
\[
2x+3=17.
\]

Donc :
\[
2x=14,
\]
et :
\[
x=7.
\]

Cette question inverse peut également être lue sur un tableau ou un graphique.

\section{7. Du tableau au graphique}

Chaque couple :
\[
(x;y)
\]
devient un point du repère.

Par exemple :
\[
(0;3),\quad(1;5),\quad(2;7).
\]

L'abscisse représente la grandeur d'entrée.

L'ordonnée représente la grandeur associée.

\section{8. Lire une valeur sur un graphique}

\coursefigure{/images/mathematiques/4eme/fonctions-lecture.svg}{Courbe avec projections pointillées permettant de lire une valeur associée à une abscisse.}{Pour une abscisse donnée, on rejoint le graphique puis on lit l'ordonnée correspondante sur l'autre axe.}

\begin{methode}
\begin{enumerate}
\item partir de la valeur connue sur un axe ;
\item rejoindre le graphique parallèlement à l'autre axe ;
\item repérer le point d'intersection ;
\item projeter sur l'autre axe ;
\item lire la valeur et son unité.
\end{enumerate}
\end{methode}

\section{9. Lecture inverse}

Si l'on connaît une ordonnée et que l'on cherche les abscisses correspondantes, on procède dans l'autre sens.

Une même ordonnée peut parfois correspondre à plusieurs abscisses.

Il faut donc observer tout le graphique.

\section{10. Exemple de température}

Une courbe donne la température au cours d'une journée.

À :
\[
14\ \text{h},
\]
on peut lire :
\[
23^\circ\text{C}.
\]

À l'inverse, une température de :
\[
18^\circ\text{C}
\]
peut être atteinte à deux moments différents, par exemple le matin et le soir.

Le graphique permet de voir cette multiplicité.

\section{11. Formule géométrique}

Pour un carré de côté $c$ :
\[
A=c^2.
\]

On peut produire :

\begin{tabular}{c|c|c|c|c}
$c$ & $0$ & $1$ & $2$ & $3$ \\
$A$ & $0$ & $1$ & $4$ & $9$ \\
\end{tabular}

Les points ne sont pas alignés sur une droite.

La dépendance n'est pas proportionnelle.

\section{12. Proportionnalité comme cas particulier}

Si :
\[
y=4x,
\]
alors la relation est proportionnelle.

Le graphique passe par :
\[
O(0;0)
\]
et les points sont alignés.

La proportionnalité est donc un cas particulier de dépendance entre grandeurs.

\section{13. Relation non proportionnelle}

Si :
\[
y=4x+3,
\]
alors :
\[
x=0\Rightarrow y=3.
\]

La représentation ne passe pas par l'origine.

La relation n'est pas proportionnelle.

\section{14. Données discrètes}

Le nombre de billets achetés peut prendre des valeurs :
\[
0,\ 1,\ 2,\ 3,\ldots
\]

Des valeurs comme :
\[
1{,}5
\]
billet n'ont pas de sens.

On représente alors plutôt des points isolés.

\begin{attention}
Relier systématiquement tous les points peut suggérer des valeurs intermédiaires qui n'existent pas dans le contexte.
\end{attention}

\section{15. Données continues}

Une durée ou une température peut prendre des valeurs intermédiaires.

Dans ce type de contexte, une courbe continue peut avoir du sens.

Le choix graphique dépend donc de la nature des grandeurs.

\section{16. Tableau incomplet}

On connaît :
\[
y=5x-2.
\]

Si :
\[
x=3,
\]
alors :
\[
y=13.
\]

Si :
\[
y=28,
\]
on résout :
\[
5x-2=28.
\]

Donc :
\[
5x=30,
\]
et :
\[
x=6.
\]

\section{17. Comparer deux tarifs}

Tarif A :
\[
A=4x.
\]

Tarif B :
\[
B=2x+12.
\]

Pour :
\[
x=3,
\]
on obtient :
\[
A=12,\qquad B=18.
\]

Pour :
\[
x=10,
\]
on obtient :
\[
A=40,\qquad B=32.
\]

Le tarif le plus avantageux dépend donc de la valeur de $x$.

\section{18. Point d'égalité de deux tarifs}

On cherche :
\[
4x=2x+12.
\]

Donc :
\[
2x=12,
\]
puis :
\[
x=6.
\]

Pour :
\[
x=6,
\]
les deux tarifs coûtent :
\[
24\ \text{€}.
\]

Cette égalité peut être repérée algébriquement ou graphiquement par l'intersection de deux représentations.

\section{19. Graphique et estimation}

Une lecture graphique fournit souvent une valeur approchée.

Si le point d'intersection semble avoir pour abscisse :
\[
6{,}1,
\]
la formule peut permettre de vérifier si la valeur exacte est plutôt :
\[
6.
\]

Les représentations se complètent.

\section{20. Produire une formule}

Un rectangle possède une largeur fixe :
\[
4\ \text{cm}
\]
et une longueur variable $x$.

Son aire vaut :
\[
A=4x.
\]

Son périmètre vaut :
\[
P=2x+8.
\]

Deux grandeurs différentes peuvent donc dépendre du même paramètre selon des règles différentes.

\section{21. Interpréter les axes}

Un graphique sans titre ni unité est incomplet.

L'axe horizontal peut représenter :
\[
\text{temps en min}
\]
et l'axe vertical :
\[
\text{distance en km}.
\]

Échanger les axes change le sens des coordonnées.

\section{22. Valeur maximale et minimale}

Sur une courbe, on peut rechercher :
\begin{itemize}
\item la plus grande valeur atteinte ;
\item la plus petite ;
\item les périodes de hausse ;
\item les périodes de baisse ;
\item les moments où une valeur seuil est franchie.
\end{itemize}

Ces lectures utilisent le graphique sans nécessairement disposer d'une formule.

\section{23. Limite du niveau 4e}

\begin{attention}
En 4e, on développe la dépendance entre grandeurs, la lecture de tableaux, formules et graphiques.

On n'exige pas encore une formalisation systématique avec toute la notation et tout le vocabulaire abstrait des fonctions qui seront approfondis ensuite.
\end{attention}

L'objectif est avant tout de savoir :
\[
\text{interpréter et changer de représentation}.
\]

\section{24. Méthode de changement de représentation}

\begin{methode}
\begin{enumerate}
\item identifier les deux grandeurs ;
\item préciser celle qui sert d'entrée ;
\item appliquer la formule ou lire le tableau ;
\item construire les couples ;
\item placer les points si un graphique est demandé ;
\item interpréter les valeurs dans le contexte ;
\item vérifier les unités.
\end{enumerate}
\end{methode}

\section{25. Erreurs fréquentes}

\begin{itemize}
\item Inverser abscisse et ordonnée.
\item Lire un graphique sans vérifier l'échelle.
\item Confondre toutes les dépendances avec des proportionnalités.
\item Relier des points lorsque les valeurs intermédiaires n'ont pas de sens.
\item Oublier les frais fixes dans une formule.
\item Inventer une formule à partir d'un seul couple de valeurs.
\item Donner une lecture graphique comme valeur exacte sans tenir compte de la précision.
\end{itemize}

\section{À retenir}

\begin{aretenir}
Une grandeur peut être donnée en fonction d'une autre.

Une même relation peut être représentée par :
\[
\boxed{\text{une formule, un tableau ou un graphique}}.
\]

Chaque couple du tableau donne un point du graphique.

Une relation de proportionnalité est un cas particulier dont le graphique est aligné avec l'origine.

En 4e, savoir passer d'une représentation à une autre et interpréter le contexte est essentiel.
\end{aretenir}
